\begin{frame}{Плани щодо дипломної роботи}
    \begin{block}{Заплановані завдання}
        \begin{itemize}
            \item \textbf{Удосконалення генерації підказок для багатоваріантності:} 
                розробка алгоритмів семантично точного розміщення колірних маркерів з урахуванням контурів та меж об'єктів.
            \item \textbf{Розширення набору базових моделей:} 
                інтеграція готових моделей інших архітектур для проведення більш вичерпного та репрезентативного порівняльного аналізу.
            \item \textbf{Комплексне користувацьке дослідження:} 
                проведення масштабного суб'єктивного оцінювання як візуальної якості зображень, згенерованих запропонованою моделлю, так і зручності взаємодії з розробленим застосунком.
        \end{itemize}
    \end{block}
\end{frame}

\begin{frame}{Перспективи подальших досліджень}
    \begin{itemize}
        \item \textbf{Теоретичне дослідження перспективних архітектур:}
            аналіз можливостей інтеграції дифузійних процесів із SSM (наприклад, Diffusion-Mamba) для задачі колоризації зображень.
        \item \textbf{Обробка відеопотоку:}
            адаптація моделі для колоризації відео із забезпеченням часової узгодженості.
        \item \textbf{Розширення функціоналу програмного комплексу:}
            впровадження додаткових інструментів для підвищення зручності користувацької взаємодії.
    \end{itemize}
\end{frame}